\documentclass[a4paper,12pt]{article}
\usepackage[T2A]{fontenc}
\usepackage[utf8]{inputenc}
\usepackage[russian]{babel}
\usepackage{amsmath,amssymb,mathtools}
\usepackage{icomma}
\usepackage{geometry}
\geometry{left=22mm,right=22mm,top=21mm,bottom=23mm}
\usepackage{microtype}
\usepackage{tikz}
\usetikzlibrary{calc,arrows.meta,positioning,shapes.geometric}
\usepackage{pgfplots}
\pgfplotsset{compat=1.18}
\usepackage{tabularx,booktabs,array}
\usepackage{enumitem}
\usepackage{fancyhdr}
\usepackage{parskip}
\usepackage{hyperref}
\usepackage{xcolor}
\definecolor{accent}{HTML}{008080}
\definecolor{darkaccent}{HTML}{005656}
\definecolor{textgray}{HTML}{454545}
\definecolor{lightgray}{HTML}{E5ECEC}
\hypersetup{colorlinks=true,linkcolor=darkaccent,urlcolor=darkaccent,pdftitle={Чек-лист: функции, графики и параметры},pdfauthor={Лёвин Артём Александрович}}
\setlength{\parindent}{0pt}
\setlength{\parskip}{5pt}
\linespread{1.13}
\setlist[itemize]{leftmargin=6mm,topsep=3pt,itemsep=2pt}
\setlist[enumerate]{leftmargin=7mm,topsep=4pt,itemsep=6pt}
\renewcommand{\arraystretch}{1.18}
\pagestyle{fancy}
\fancyhf{}
\fancyhead[L]{\footnotesize\color{textgray}Функции, графики и параметры}
\fancyhead[R]{\footnotesize\color{textgray}Анна Трапезникова}
\fancyfoot[C]{\footnotesize\color{textgray}\thepage}
\renewcommand{\headrulewidth}{0.25pt}
\setlength{\headheight}{15pt}
\newcommand{\sect}[1]{\vspace{4pt}{\large\bfseries\color{darkaccent}#1}\par\vspace{2pt}}
\newcommand{\idea}[1]{\textcolor{darkaccent}{\textbf{Ключевая идея.}} #1}
\newcommand{\warn}[1]{\textcolor{darkaccent}{\textbf{Важно.}} #1}
\newcommand{\ex}[1]{\textbf{Пример.} #1}
\newcommand{\chk}{\(\square\)\,}
\begin{document}
\thispagestyle{empty}
\begin{center}
\vspace*{25mm}
{\large\color{darkaccent}МАТЕМАТИКА}\par
\vspace{19mm}
{\LARGE\color{textgray}Чек-лист}\par
\vspace{12mm}
{\Large\color{accent}Функции, графики и параметры}\par
\vspace{3mm}
{\large\color{accent}ОДЗ, разложение на множители, модуль}\par
\vspace{17mm}
{\normalsize\color{textgray}Алгебра и графические методы решения задач}\par
\vspace{37mm}
{\large Анна Трапезникова}\par
\vspace{13mm}
{\normalsize 10 ноября 2026 г.}\par
\vfill
{\small Лёвин Артём Александрович}\par
{\small эксклюзивно для Анны Трапезниковой}\par
\vspace{21mm}
\begin{tikzpicture}[x=1cm,y=1cm]
\draw[accent,line width=0.9pt] (0,0) .. controls (1.2,0.55) and (2.6,-0.55) .. (4,0) .. controls (5.2,0.5) and (6.3,-0.55) .. (8,0);
\fill[accent] (4,0) circle (1.3pt);
\end{tikzpicture}
\end{center}
\clearpage

\sect{1. Что нужно уметь после занятия}
\begin{itemize}
\item \chk Выбирать нужную формулу для каждого промежутка кусочно заданной функции и учитывать включённые границы.
\item \chk Находить \textbf{область допустимых значений} (ОДЗ) до упрощения дроби.
\item \chk Разлагать многочлен на множители заменой \(t=x^2\), формулами сокращённого умножения и схемой Горнера.
\item \chk Строить графики, оставляя \textbf{выколотые точки} в местах, запрещённых исходной функцией.
\item \chk Отличать \(y=|f(x)|\), \(y=f(|x|)\) и функции с модулем лишь в части выражения.
\item \chk Определять число пересечений графика с \(y=c\) и \(y=kx\).
\end{itemize}

\sect{2. Как вычислять значения кусочной функции}
Если функция задана несколькими формулами, сначала найдите промежуток, которому принадлежит заданный \(x\), и только затем подставляйте его в соответствующую формулу.
\[
 f(x)=\begin{cases}x+2,&x<1,\\5-x,&x\geqslant 1.\end{cases}
 \qquad f(-1)=1,\quad f(1)=4.
\]
Точка \(x=1\) относится ко \textbf{второй} формуле: её условие содержит знак равенства. На графике закрашенная точка означает включённое значение, незакрашенная --- исключённое.

\sect{3. ОДЗ: почему нельзя терять ограничения}
Для дроби \(\frac{P(x)}{Q(x)}\) прежде всего запишите \(Q(x)\ne0\). Сокращение множителей упрощает запись функции, но \textbf{не возвращает} запрещённые значения \(x\).

\ex{Исследуем функцию}
\[
F(x)=\frac{x^4-25x^2+144}{(x-4)(x+3)}.
\]
\textbf{Шаг 1.} По знаменателю \(x\ne4\), \(x\ne-3\).

\textbf{Шаг 2.} В числителе делаем замену \(t=x^2\):
\[
t^2-25t+144=(t-16)(t-9).
\]
Восстанавливаем \(x\), затем используем разность квадратов:
\[
\begin{aligned}
x^4-25x^2+144
 &=(x^2-16)(x^2-9)\\
 &=(x-4)(x+4)(x-3)(x+3).
\end{aligned}
\]
\textbf{Шаг 3.} Сокращаем только при уже записанном ОДЗ:
\[
F(x)=(x+4)(x-3)=x^2+x-12,\quad x\ne4,-3.
\]
\idea{График совпадает с параболой \(y=x^2+x-12\), из которой удалены точки \((4;8)\) и \((-3;-6)\).}

\clearpage
\sect{4. Параметр \(c\): горизонтальная прямая}
Прямая \(y=c\) горизонтальна. Число её общих точек с графиком функции равно числу \textbf{допустимых} решений уравнения \(F(x)=c\).

Для параболы из предыдущего примера:
\[
y=x^2+x-12=\left(x+\frac12\right)^2-\frac{49}{4}.
\]
Вершина имеет координаты \(\left(-\frac12;-\frac{49}{4}\right)\). На уровне вершины горизонтальная прямая касается параболы в одной точке. Ещё две прямые проходят через выколотые точки, оставляя ровно одно пересечение.

\begin{center}
\begin{tikzpicture}
\begin{axis}[
 width=12.6cm,height=8.4cm, axis lines=middle,
 xmin=-6,xmax=5.3,ymin=-14,ymax=13,
 xtick={-5,-3,-1,0,2,4},ytick={-12,-6,0,8},
 xlabel={$x$},ylabel={$y$},
 xlabel style={at={(axis description cs:0.98,0.45)},anchor=west},
 ylabel style={at={(axis description cs:0.48,0.98)},anchor=south},
 ticklabel style={font=\small},clip=false,
 ]
\addplot[accent,very thick,domain=-5.5:4.8,samples=160] {x^2+x-12};
\addplot[only marks,mark=*,mark size=2.2pt,darkaccent] coordinates {(-0.5,-12.25)};
\addplot[only marks,mark=o,mark size=3pt,mark options={line width=1.2pt,draw=accent,fill=white}] coordinates {(-3,-6) (4,8)};
\addplot[dashed,gray,domain=-6:5.2] {-12.25};
\addplot[dashed,gray,domain=-6:5.2] {-6};
\addplot[dashed,gray,domain=-6:5.2] {8};
\node[anchor=south west,font=\small] at (axis cs:4,8) {$B(4;8)$};
\node[anchor=north east,font=\small] at (axis cs:-3,-6) {$A(-3;-6)$};
\node[anchor=north east,font=\small] at (axis cs:-0.5,-12.25) {$V$};
\end{axis}
\end{tikzpicture}
\end{center}
\textbf{Ровно одно пересечение} получится при
\[
\boxed{c\in\left\{-\frac{49}{4},\,-6,\,8\right\}.}
\]
\warn{Уровень выколотой точки может оставить одно пересечение только тогда, когда на той же высоте существует другая допустимая точка. Проверяйте это для каждого графика отдельно.}

\sect{5. Схема Горнера: полный расклад многочлена}
Для многочлена \(P(x)=x^4+x^3-7x^2-x+6\) целые корни ищем среди делителей свободного члена \(6\): \(\pm1,\pm2,\pm3,\pm6\).

Поскольку \(P(1)=0\), делим на \(x-1\) по схеме Горнера. Коэффициенты \(1,1,-7,-1,6\) дают остаток \(0\) и частное
\[
x^3+2x^2-5x-6.
\]
Для него проверяем \(x=-1\): этот корень даёт множитель \(x+1\), а оставшееся частное равно \(x^2+x-6=(x+3)(x-2)\). Поэтому
\[
\boxed{P(x)=(x-1)(x+1)(x+3)(x-2).}
\]
\warn{Корню \(r\) соответствует множитель \(x-r\). При корне \(-1\) множителем будет \(x+1\). Проверка: произведение даёт исходный многочлен.}

\clearpage
\sect{6. Модуль и графические преобразования}
\textbf{Три внешне похожие записи требуют разных действий.}
\begin{center}
\small
\begin{tabularx}{\linewidth}{@{}p{31mm}X@{}}
\toprule
Форма & Как строить график \\
\midrule
\(y=|f(x)|\) & Постройте \(y=f(x)\). Части ниже оси \(Ox\) отразите \textbf{вверх} относительно \(Ox\). \\
\(y=f(|x|)\) & Оставьте правую половину графика \(y=f(x)\) и отразите её относительно оси \(Oy\). \\
\(y=\frac1{|x|+1}\) & Модуль стоит внутри выражения: раскройте \(|x|\) при \(x\ge0\) и при \(x<0\), затем постройте обе части. \\
\bottomrule
\end{tabularx}
\end{center}
Для последнего примера:
\[
\frac1{|x|+1}=\begin{cases}
\frac1{x+1},&x\ge0,\\[2pt]
\frac1{1-x},&x<0.
\end{cases}
\]
\ex{Из гиперболы \(y=1/x\) получаем \(y=1/|x|\), отражая её правую ветвь относительно \(Oy\). Затем знак минус перед дробью отражает график относительно \(Ox\): \(y=-1/|x|\).}

\sect{7. Сокращение выражения с модулем}
Возьмём дробно-рациональную функцию
\[
G(x)=\frac{3|x|-1}{|x|-3x^2}.
\]
Используем \(x^2=|x|^2\). Её знаменатель равен \(|x|(1-3|x|)\); поэтому ОДЗ:
\[
|x|\ne0,\quad 1-3|x|\ne0
\quad\Longleftrightarrow\quad x\ne0,\ \ x\ne\pm\frac13.
\]
После вынесения минуса в числителе
\[
G(x)=\frac{-(1-3|x|)}{|x|(1-3|x|)}
=-\frac1{|x|},\quad x\ne0,\pm\frac13.
\]
В графике \(y=-1/|x|\) отсутствуют точки \(\left(-\frac13;-3\right)\) и \(\left(\frac13;-3\right)\). При \(x=0\) имеется вертикальная асимптота, поэтому конечной выколотой точки там нет.

\begin{center}
\begin{tikzpicture}
\begin{axis}[
 width=11.9cm,height=6.7cm,axis lines=middle,
 xmin=-3.2,xmax=3.2,ymin=-5.2,ymax=0.8,
 xtick={-2,-1,0,1,2},ytick={-4,-3,-2,-1,0},
 xlabel={$x$},ylabel={$y$},ticklabel style={font=\small},
 xlabel style={at={(axis description cs:0.98,0.87)},anchor=west},
 ylabel style={at={(axis description cs:0.5,1)},anchor=south},clip=false]
\addplot[accent,very thick,domain=-3.1:-0.205,samples=125] {1/x};
\addplot[accent,very thick,domain=0.205:3.1,samples=125] {-1/x};
\addplot[only marks,mark=o,mark size=3pt,mark options={draw=accent,fill=white,line width=1pt}] coordinates {(-0.333333,-3) (0.333333,-3)};
\node[anchor=north east,font=\scriptsize] at (axis cs:-0.333333,-3) {$x=-\frac13$};
\node[anchor=north west,font=\scriptsize] at (axis cs:0.333333,-3) {$x=\frac13$};
\end{axis}
\end{tikzpicture}
\end{center}

\clearpage
\sect{8. Параметр \(k\): прямая через начало координат}
Прямая \(y=kx\) всегда проходит через \((0;0)\). Число \(k\) задаёт её угловой коэффициент; при \(k=0\) это ось \(Ox\).

\ex{Пусть \(H(x)=|x|+2\), причём \(x\ne-2\) и \(x\ne2\). Найдём все \(k\), при которых \(y=kx\) и график \(H\) не пересекаются.}
\[
H(x)=\begin{cases}x+2,&x\ge0,\\-x+2,&x<0.\end{cases}
\]
\textbf{Правая ветвь.} Решаем \(kx=x+2\):
\[
(k-1)x=2,\qquad x=\frac2{k-1}.
\]
Это пересечение существует при \(k>1\), но если \(k=2\), получаем запрещённое \(x=2\): пересечение исчезает.

\textbf{Левая ветвь.} Решаем \(kx=-x+2\):
\[
(k+1)x=2,\qquad x=\frac2{k+1}.
\]
Для левой ветви нужно \(x<0\), то есть \(k<-1\). При \(k=-2\) решение \(x=-2\) запрещено.

\textbf{Общий вывод.} При всех \(-1\le k\le1\) пересечений нет: это следует из направлений лучей графика. Кроме того, подходят два исключительных наклона, проходящих через выколотые точки:
\[
\boxed{k\in[-1;1]\cup\{-2;2\}.}
\]
\warn{Параллельность прямой отдельной ветви даёт только граничные случаи \(k=\pm1\). Между ними также нет пересечений: обязательно проверяйте весь интервал значений параметра.}

\begin{center}
\begin{tikzpicture}
\begin{axis}[
 width=11.7cm,height=6.7cm,axis lines=middle,
 xmin=-4.4,xmax=4.4,ymin=-1,ymax=7,
 xtick={-4,-2,0,2,4},ytick={0,2,4,6},
 xlabel={$x$},ylabel={$y$},ticklabel style={font=\small},clip=false,
 xlabel style={at={(axis description cs:0.98,0.1)},anchor=west},
 ylabel style={at={(axis description cs:0.5,0.99)},anchor=south}]
\addplot[accent,very thick,domain=-4.2:0] {-x+2};
\addplot[accent,very thick,domain=0:4.2] {x+2};
\addplot[only marks,mark=o,mark size=3pt,mark options={draw=accent,fill=white,line width=1pt}] coordinates {(-2,4) (2,4)};
\addplot[gray,dashed,domain=-0.45:3.25] {x};
\addplot[gray,dashed,domain=-3.25:0.45] {-x};
\addplot[darkaccent,dashdotted,domain=-0.3:3.25] {2*x};
\addplot[darkaccent,dashdotted,domain=-3.25:0.3] {-2*x};
\end{axis}
\end{tikzpicture}
\end{center}

\clearpage
\begin{center}
{\LARGE\color{darkaccent}Алгоритм: исследование графика дробной функции}\par
\vspace{15mm}
\begin{tikzpicture}[
 x=1cm,y=1cm, >=Stealth,
 boxstep/.style={draw=darkaccent,rounded corners=2pt,align=center,text width=9.3cm,minimum height=0.95cm,font=\small,inner sep=4pt},
 term/.style={ellipse,draw=accent,align=center,minimum width=4.6cm,minimum height=0.80cm,font=\small},
 ask/.style={diamond,draw=darkaccent,aspect=2.6,align=center,text width=2.8cm,inner sep=2pt,font=\small},
 route/.style={-Stealth,thick,draw=darkaccent}
 ]
\node[term] (start) at (0,0) {Дана функция $F(x)$};
\node[boxstep] (odz) at (0,-1.35) {Найти ОДЗ исходного выражения};
\node[boxstep] (simp) at (0,-2.70) {Разложить на множители и упростить формулу, сохраняя ОДЗ};
\node[ask] (mod) at (0,-4.15) {Есть модуль?};
\node[boxstep,text width=5.15cm] (yes) at (-3.35,-5.85) {Раскрыть модуль по случаям или выполнить допустимое отражение};
\node[boxstep,text width=5.15cm] (no) at (3.35,-5.85) {Построить график упрощённой формулы};
\node[boxstep] (plot) at (0,-7.55) {Объединить полученные ветви на одной координатной плоскости};
\node[boxstep] (holes) at (0,-8.90) {Удалить точки с запрещёнными $x$ и отметить асимптоты};
\node[ask] (kind) at (0,-10.38) {Какая прямая?};
\node[boxstep,text width=5.15cm] (hor) at (-3.35,-12.08) {$y=c$: проверить горизонтальные пересечения};
\node[boxstep,text width=5.15cm] (rot) at (3.35,-12.08) {$y=kx$: решить $F(x)=kx$, проверить ветви и ОДЗ};
\node[term] (end) at (0,-13.75) {Записать полный набор параметров};
\draw[route] (start)--(odz);
\draw[route] (odz)--(simp);
\draw[route] (simp)--(mod);
\draw[route] (mod.south west) -| node[pos=0.15,above left,font=\scriptsize]{да} (yes.north);
\draw[route] (mod.south east) -| node[pos=0.15,above right,font=\scriptsize]{нет} (no.north);
\draw[draw=darkaccent,thick] (yes.south) -- (-3.35,-6.72) -- (0,-6.72);
\draw[draw=darkaccent,thick] (no.south) -- (3.35,-6.72) -- (0,-6.72);
\draw[route] (0,-6.72) -- (plot.north);
\draw[route] (plot)--(holes);
\draw[route] (holes)--(kind);
\draw[route] (kind.south west) -| node[pos=0.15,above left,font=\scriptsize]{$y=c$} (hor.north);
\draw[route] (kind.south east) -| node[pos=0.15,above right,font=\scriptsize]{$y=kx$} (rot.north);
\draw[draw=darkaccent,thick] (hor.south) -- (-3.35,-12.97) -- (0,-12.97);
\draw[draw=darkaccent,thick] (rot.south) -- (3.35,-12.97) -- (0,-12.97);
\draw[route] (0,-12.97) -- (end.north);
\end{tikzpicture}
\end{center}

\clearpage
\sect{9. Типичные ошибки и самопроверка}
\begin{itemize}
\item \textbf{Сокращение без ОДЗ.} Запрещённые значения сохраняются после сокращения.
\item \textbf{Неправильный знак множителя.} Корень \(-a\) даёт множитель \(x+a\).
\item \textbf{Не тот тип отражения.} Для \(|f(x)|\) отражаем относительно \(Ox\), для \(f(|x|)\) --- относительно \(Oy\).
\item \textbf{Выколотая точка нарисована закрашенной.} Сначала вычислите координаты по упрощённой формуле, затем удалите точку.
\item \textbf{Параметр задан отдельными значениями вместо интервала.} Проверяйте все возможные положения прямой между граничными наклонами.
\item \textbf{Перепутаны \(c\) и \(k\).} Прямая \(y=c\) перемещается вертикально, прямая \(y=kx\) меняет наклон вокруг начала координат.
\end{itemize}

\idea{При повторном решении пройдите задачу самостоятельно от начала до конца: ОДЗ \(\to\) преобразование \(\to\) график \(\to\) удалённые точки \(\to\) проверка параметра.}

\clearpage
\begin{center}{\LARGE\color{darkaccent}Тренировочный блок}\par\vspace{2mm}{\small Путь А. 10 заданий от базовых к более содержательным.}\end{center}
\vspace{6mm}
\begin{enumerate}
\item Для \(f(x)=\begin{cases}x+2,&x<1,\\5-x,&x\ge1\end{cases}\) найдите \(f(-1), f(0), f(1), f(3)\).
\item Разложите на линейные множители \(x^4-13x^2+36\).
\item Разложите по схеме Горнера (или другим школьным способом) \(x^4+x^3-7x^2-x+6\).
\item Упростите дробь \(\displaystyle\frac{x^4-25x^2+144}{(x-4)(x+3)}\). Укажите ОДЗ и координаты выколотых точек.
\item Для функции из задания 4 найдите все \(c\), при которых прямая \(y=c\) имеет с графиком ровно одну общую точку.
\item Постройте \(y=|x^2-4|\); найдите нули функции и значение при \(x=0\).
\item Упростите \(\displaystyle\frac{3|x|-1}{|x|-3x^2}\). Укажите все исключённые значения \(x\) и координаты выколотых точек.
\item Раскройте модуль по определению для \(\displaystyle y=\frac1{|x|+1}\). Укажите область определения и значения при \(x=0,1,-1\).
\item График \(y=|x|+2\) задан при \(x\ne-2,2\). Найдите все \(k\), для которых прямая \(y=kx\) не имеет с ним общих точек.
\item Исследуйте \(\displaystyle y=\frac{x^4-10x^2+9}{(x-1)(x+3)}\). Найдите ОДЗ, упрощённую формулу и все \(c\), для которых прямая \(y=c\) имеет с графиком ровно одну общую точку.
\end{enumerate}

\vfill
\textcolor{textgray}{\small Упражнения 4--5 и 7--10 рекомендуется выполнять с аккуратным графиком и отдельной проверкой исключённых значений.}

\clearpage
\begin{center}{\LARGE\color{darkaccent}Ответы к тренировочному блоку}\end{center}
\vspace{8mm}
\small
\begin{tabularx}{\linewidth}{@{}p{11mm}X@{}}
\toprule
№ & Ответ \\
\midrule
1 & \(1;\ 2;\ 4;\ 2\). \\
2 & \((x-2)(x+2)(x-3)(x+3)\). \\
3 & \((x-1)(x+1)(x-2)(x+3)\). \\
4 & \(x^2+x-12\), \(x\ne4,-3\); выколотые точки \((4;8)\), \((-3;-6)\). \\
5 & \(c\in\left\{-\dfrac{49}{4},-6,8\right\}\). \\
6 & Нули: \(x=-2,2\). Значение \(y(0)=4\); части параболы ниже \(Ox\) отражены вверх. \\
7 & \(-\dfrac1{|x|}\), \(x\ne0,\pm\dfrac13\); выколотые точки \(\left(-\dfrac13;-3\right)\), \(\left(\dfrac13;-3\right)\). \\
8 & \(y=\begin{cases}\dfrac1{x+1},&x\ge0,\\[1pt]\dfrac1{1-x},&x<0.\end{cases}\) \(D=\mathbb R\); \(y(0)=1,\ y(1)=y(-1)=\dfrac12\). \\
9 & \(k\in[-1;1]\cup\{-2;2\}\). \\
10 & \(x\ne1,-3\); \(y=x^2-2x-3\); \(c=12\). При \(c=-4\) вершина \((1;-4)\) выколота, пересечений нет. \\
\bottomrule
\end{tabularx}
\normalsize
\vspace{16mm}
\sect{Итоговая памятка}
\small
\begin{tabularx}{\linewidth}{@{}p{43mm}X@{}}
\toprule
Сигнал в задаче & Полезное действие \\
\midrule
Есть знаменатель & Найдите ОДЗ до сокращения. \\
Четвёртая и вторая степени & Проверьте замену \(t=x^2\). \\
Многочлен высокой степени & Ищите рациональный корень, используйте схему Горнера. \\
Модуль в формуле & Определите его область действия и нужный способ построения. \\
Прямая \(y=c\) & Считайте горизонтальные пересечения. \\
Прямая \(y=kx\) & Исследуйте наклон и проверяйте допустимость точек. \\
После сокращения & Верните на график информацию об исключённых \(x\). \\
\bottomrule
\end{tabularx}
\end{document}
